\subsection{\normalsize Compresión de memoria con pérdida de precisión}
Dado que la representación exacta de estados cuánticos es ineficiente en términos
de memoria, se han explorado técnicas de compresión con pérdida de precisión. Estos
enfoques sacrifican una pequeña cantidad de fidelidad numérica a cambio de una
reducción significativa en el uso de memoria, permitiendo la simulación de
circuitos más grandes.

Uno de los primeros enfoques en este ámbito fue el de \textcite{wu2018amplitude, wu2019full},
quienes introdujeron la compresión adaptativa con error acotado. Utilizando la
biblioteca SZ, reportaron reducciones del tamaño del vector de estado de hasta en un factor
de 16 y fidelidades superiores al 99.9~\% en los casos experimentales estudiados. Su principal ventaja es
que permite ampliar el número de qubits simulados sin un incremento
proporcional en los requisitos de memoria. Sin embargo, la compresión y
descompresión recurrente introduce sobrecarga computacional, afectando la eficiencia
temporal de la simulación.

En un esfuerzo por maximizar la escala de la simulación cuántica de estado
completo, \textcite{Zhang2025BMQSim} presentaron \emph{BMQSim}, un marco que
extiende \emph{SV-Sim} mediante el uso de \textbf{compresores con pérdida de
precisión y control de error punto a punto} ejecutados directamente en GPU (por
ejemplo, variantes de la familia SZ/cuSZ). Su diseño
combina:
\begin{itemize}
    \item Una estrategia de \textit{circuit partitioning} que disminuye la
        frecuencia de (des)compresión.

    \item Un \textit{pipeline} solapado que oculta en gran medida el coste de
        mover y comprimir datos.

    \item Una gestión de memoria jerárquica (VRAM + SSD mediante GPUDirect~Storage)
        que asigna dinámicamente espacio a los bloques comprimidos.
\end{itemize}
Con estas técnicas, BMQSim reduce el consumo de memoria en más de un orden de magnitud
y mantiene fidelidades superiores al 99~\% en los experimentos reportados sin penalizar de forma apreciable el
tiempo de simulación. Su eficacia, sin embargo, está condicionada a la
disponibilidad de hardware acelerado con GPUs de alto rendimiento y
unidades NVMe rápidas, lo que puede restringir su adopción en plataformas de menor
capacidad computacional.

Más recientemente, \textcite{Huffman2024} exploraron la cuantización de
amplitudes como un mecanismo para reducir la precisión numérica sin afectar la
fidelidad global del sistema. Demostraron que, bajo las configuraciones evaluadas,
representaciones con tan solo 7 bits por amplitud permiten mantener una fidelidad
superior al 99~\%, lo que sugiere que
la precisión completa de punto flotante puede ser innecesaria en muchas simulaciones.
Aunque esta técnica ofrece un enfoque simple y eficiente para reducir el uso
de memoria, su aplicabilidad a circuitos de gran escala aún requiere validación
empírica.

Dentro de esta familia de técnicas, ZFP constituye una referencia relevante para
arreglos numéricos de punto flotante. Su esquema divide el arreglo en bloques
pequeños e independientes y transforma los valores de cada bloque a una
representación más compacta antes de codificar sus coeficientes según un
criterio de tasa, precisión o error tolerado \cite{Lindstrom2014ZFP}. Esta
organización permite comprimir y descomprimir fragmentos del arreglo sin
depender del resto, una propiedad útil cuando el acceso a los datos es parcial o
segmentado. A diferencia de los enfoques sin pérdida, ZFP no busca preservar
igualdad exacta bit a bit, sino controlar el compromiso entre reducción de
memoria y desviación numérica. Por esta razón, en este trabajo se toma como
referencia de compresión con pérdida controlada frente a la estrategia exacta
basada en Blosc.

\textcite{DiazToro2025} propone un esquema de \textbf{vector de estado
comprimido} que emplea compresores con pérdida de precisión, principalmente
\emph{ZFP} en modo \emph{fixed-rate}, complementado por \emph{SZ}, para reducir
la huella de memoria del simulador. El autor reporta reducciones del \textbf{10--20~\%}
en casos base y, en configuraciones de mayor escala, ahorros que alcanzan hasta
un \textbf{75~\%} de la memoria requerida. \emph{ZFP} opera sobre bloques independientes
de $4^{d}$ valores, de modo que la (des)compresión de cada bloque se realiza en
completo aislamiento del resto; esto permite solapar dichas operaciones con el
cómputo principal y amortiguar la sobrecarga introducida. Aunque la compresión
incrementa el tiempo de ejecución, el impacto se compensa al transmitir los
fragmentos del vector en formato comprimido dentro de una arquitectura
distribuida basada en \emph{MPI}, reduciendo el volumen de comunicación y
habilitando simulaciones de hasta $30$ qubits en los experimentos reportados. Finalmente,
el estudio concluye que esta estrategia de \textit{estado completo comprimido
distribuido} resulta más fiable y escalable que la \textit{poda de estados de
baja probabilidad}, la cual presenta descensos de fidelidad y sobrecarga adicional
que la hacen desaconsejable.
